\section{Properties of the Phase Readout}
\label{app:proofs}

\myparagraph{Proof of Proposition~\ref{prop:exact}.}
\begin{proof}
Set $\omega = e^{2\pi i/k}$, $z_0=z_1=1$ initially, $\theta_0(o)=0$,
$\theta_1(o)=\tfrac{2\pi}{k}\,e(o)$ and $\mathrm{inj}\equiv 0$, with $e(o)$ the
event bit encoded in $o$. Then $z_0=1$ at every step and $z_1=\omega^{c}$ with
$c=\sum_t e_t$. Readout weights $w_{a0}=1$ and $w_{a1}=\omega^{-a}$ give
\begin{equation}
\psi_a = 1 + \omega^{\,c-a},\qquad
|\psi_a|^2 = 2 + 2\cos\!\Big(\tfrac{2\pi}{k}(c-a)\Big),
\label{eq:interference}
\end{equation}
maximized exactly at $a\equiv c \pmod{k}$. Normalizing the state multiplies the
two coordinates by the same positive scalar and leaves the maximizing action
unchanged.
\end{proof}


\begin{proposition}[Global-phase invariance of squared-amplitude scores]
\label{prop:gauge}
With the complex-linear convention $\psi_a=\sum_jw_{aj}z_j$, the score
$|\psi_a|^2$ is unchanged by $z\mapsto e^{i\varphi}z$.
No nonzero real-linear scalar score on $(\operatorname{Re}z,\operatorname{Im}z)$ has this
invariance, and the minimum degree of a nonconstant invariant real polynomial in these
coordinates is two.
\end{proposition}
\begin{proof}
Multiplying $z$ by $e^{i\varphi}$ multiplies $\psi_a$ by the same unit-modulus scalar.
Invariance under $\varphi=\pi$ forces any real-linear invariant to be zero, while
$\operatorname{Re}(\overline{z_i}z_j)$ gives a nonconstant quadratic invariant.
\end{proof}
Relative-phase cosine decoders and Hermitian quadratic forms also have this
symmetry, as do some nonpolynomial functions.
A fixed reference channel also permits exact real-linear decoding of the phase channel,
so absolute-phase controls compare information access as well as readout form.

\begin{lemma}[Phase blindness of magnitude-only readouts]
\label{lem:blind}
In the solution of Proposition~\ref{prop:exact}, the unit-modulus coordinate
$z_1=\omega^{c}$ carries the count and $|z_0|=|z_1|=1$. Any readout depending on the
coordinate magnitudes alone, $\pi(a)=f_a(|z_0|,|z_1|)$, is independent of $c$. Its accuracy
cannot exceed that of the best constant predictor, $\max_a \Pr[a^\star=a]$. This bound
equals the chance value $1/k$ exactly for a uniform residue, as in parity.
\end{lemma}

\begin{proof}
Along the solution $|z_0|=|z_1|=1$ for every input and every $c$, so $f_a(|z_0|,|z_1|)$ does not
depend on $c$ and induces a fixed action distribution $q(a)$. Its expected accuracy is
$\sum_a q(a)\,\Pr[a^\star=a]\le \max_a \Pr[a^\star=a]$. For independent uniform bits and $k=2$ the
residue is uniform, so the bound is the chance value $1/2$; more generally the bound equals $1/k$ precisely for a uniform residue distribution.
\end{proof}

Equation~\eqref{eq:interference} is a correlation with the first Fourier
character of the cycle. Using all $k$ characters would give the squared
magnitudes of the discrete Fourier transform.

\begin{proposition}[Zero set of the phase regularizer]
\label{prop:reg}
The regularizer $R(\theta)=\big\langle\sin^2\!\big(\tfrac{k}{2}\,\theta(o)\big)\big\rangle\ge 0$
(with $\langle\cdot\rangle$ the average over observations) satisfies $R(\theta)=0$ if and only if
$\theta(o)\in\tfrac{2\pi}{k}\Z$ for every observation in a finite average (or almost surely for an expectation); its zero set consists exactly of the
per-event rotations $\theta=2\pi j/k$ for $j\in\{0,\dots,k-1\}$. Among these, every $j$ with
$\gcd(j,k)=1$ is a faithful generator that yields an exact decoder by Proposition~\ref{prop:exact}
after relabeling actions $a\mapsto ja\bmod k$, while $j=0$ is a constant phase that cannot represent
counting. The canonical choice is the generator $j=1$, that is $\theta=2\pi/k$.
\end{proposition}

\begin{proof}
$\sin^2(\tfrac{k}{2}\theta)=0$ if and only if $\tfrac{k}{2}\theta\in\pi\Z$, that is
$\theta\in\tfrac{2\pi}{k}\Z$, with representatives $\theta=2\pi j/k$ in $[0,2\pi)$ for
$j=0,\dots,k-1$. Under such a rotation the phase channel after $c$ events is $\omega^{\,jc}$. The
map $c\bmod k\mapsto jc\bmod k$ is a bijection if and only if $\gcd(j,k)=1$, and the readout of
Proposition~\ref{prop:exact} then decodes $c$ after the stated relabeling. For $j=0$ the phase
is constant and phase-blind by Lemma~\ref{lem:blind}.
\end{proof}

The degenerate zero at $j=0$ explains an empirical sensitivity: the coefficient $\lambda$ must be
small enough that the reward objective keeps the rotation away from the no-rotation basin.

\section{Extrapolation across group orders}
\label{app:kboundary}
Table~\ref{tab:kscaling} reports reward-trained extrapolation after a
curriculum of lengths 4 through 14. The decoder margin below quantifies how
closely the ideal scores are spaced as $k$ increases. Controls with exact angles
are reported in Appendix~\ref{app:additional}.

\begin{lemma}[Interference margin and its decay in $k$]
\label{lem:margin}
In the exact solution of Proposition~\ref{prop:exact}, the squared-amplitude
score of the correct residue exceeds that of every other action by
\begin{equation}
\begin{aligned}
\Delta_k &\;=\; |\psi_{a^\star}|^2 - \max_{a\neq a^\star}|\psi_a|^2\\
&\;=\; 4 - \Big(2+2\cos\tfrac{2\pi}{k}\Big) \;=\; 4\sin^2\!\tfrac{\pi}{k}.
\end{aligned}
\label{eq:margin}
\end{equation}
The margin decays as $4\pi^2/k^2$ for large $k$.
\end{lemma}

\begin{proof}
From Equation~\eqref{eq:interference}, $|\psi_a|^2=2+2\cos(\tfrac{2\pi}{k}(c-a))$, maximal at $a=c$
with value $4$; the nearest competitors $a=c\pm1$ give $2+2\cos\tfrac{2\pi}{k}$. Subtracting and
using $1-\cos\theta=2\sin^2(\theta/2)$ gives
$\Delta_k=2\big(1-\cos\tfrac{2\pi}{k}\big)=4\sin^2\tfrac{\pi}{k}$; the small-angle expansion yields
the $O(1/k^2)$ decay.
\end{proof}


\begin{table}[tbp]
\centering
\caption{Zero-shot greedy accuracy at three evaluation lengths (five seeds; training curriculum
$L\in\{4,6,\dots,14\}$). Bold marks the higher accuracy of the two models.}
\label{tab:kscaling}
\small
\setlength{\tabcolsep}{3.5pt}
\begin{tabular}{cc ccc ccc}
\toprule
 & & \multicolumn{3}{c}{\textbf{Phase rotation}} & \multicolumn{3}{c}{LSTM} \\
\cmidrule(lr){3-5} \cmidrule(lr){6-8}
$k$ & chance & L10 & L40 & L80 & L10 & L40 & L80 \\
\midrule
\wsrow $2$ & 0.50 & \textbf{1.00} & \textbf{1.00} & \textbf{0.99} & 0.50 & 0.50 & 0.50 \\
\wsrow $3$ & 0.33 & \textbf{0.91} & \textbf{0.90} & \textbf{0.70} & 0.47 & 0.43 & 0.43 \\
$4$ & 0.25 & 0.80 & 0.55 & 0.54 & \textbf{0.88} & \textbf{0.85} & \textbf{0.84} \\
$5$ & 0.20 & 0.82 & 0.57 & 0.57 & \textbf{0.84} & \textbf{0.84} & \textbf{0.84} \\
\bottomrule
\end{tabular}
\end{table}


\section{Contraction, Margin, and Stochastic Decoding}
\label{app:contractive}
Norm preservation avoids one source of lost separability. The following result makes
this precise for uniformly contractive memories.

\begin{proposition}[Contraction and robust decoding]
\label{prop:contractive}
Suppose the state update is uniformly $\gamma$-contractive in its state, $0\le\gamma<1$.
If two prefixes produce states at distance at most $D$, appending a common suffix of
length $T$ leaves their states at distance at most $D\gamma^T$. If their target labels
differ and every decoder score is $K$-Lipschitz, the two predictions cannot both have a
correct-label margin $m>0$ once $KD\gamma^T<m$. For a balanced family of $k$ prefixes
carrying all residues, with state diameter at most $D$ and a common suffix, a stochastic
decoder with $K$-Lipschitz probability coordinates has mean success at most
$1/k+KD\gamma^T$.
\end{proposition}

\begin{proof}
Repeated contraction gives $\|h_T-h_T'\|\leq D\gamma^T$. If the two
correct labels are $a$ and $b$, both predictions having margin $m$ would
require $(s_a-s_b)(h_T)-(s_a-s_b)(h_T')\geq2m$. The difference of the
scores is $2K$-Lipschitz, so this requires $m\leq KD\gamma^T$.

For the balanced family, compare each terminal state $h_T^{(r)}$ to the
reference $h_T^{(0)}$. The correct labels $a_r$ are a permutation of the
$k$ residues, so the mean sampled success is bounded by
\[
\frac{1}{k}\sum_r p_{a_r}(h_T^{(r)})
\leq \frac{1}{k}\sum_r p_{a_r}(h_T^{(0)})+KD\gamma^T
= \frac{1}{k}+KD\gamma^T.
\]
\end{proof}
A contractive update can still decode with a vanishing margin. For example,
$h_0=1$ and $h_t=\gamma(-1)^{e_t}h_{t-1}$ recover parity by
$\arg\max(h_T,-h_T)$ at every finite length in exact arithmetic.

\section{The Diagonal Complex State-Space Baseline}
\label{app:lru}
We compare linear recurrent units (LRUs) on length-10 parity over five
seeds using the same width of 16, trainer, and budget as Table~\ref{tab:mech}.
Following \citet{orvieto2023lru}, the recurrence is
$x_t=\lambda\odot x_{t-1}+\gamma\odot(Bu_t)$, with
$\lambda_j=\exp(-\exp(\nu_j)+i\exp(\phi_j))$ and
$\gamma_j=\sqrt{1-|\lambda_j|^2}$. We initialize $|\lambda_j|$ uniformly
in $[0.9,0.999]$ and its angle in $[0,2\pi)$. In the unit-modulus
controls we set $|\lambda_j|=\gamma_j=1$. The last row of
Table~\ref{tab:lru} uses the same unit-modulus recurrence and squared-amplitude readout
with the angle chosen from the observation. It trains without the
stabilizers or the length curriculum.

For a single fixed linear layer, the state is affine in the input and the
difference between two quadratic scores has degree at most two. Parity
for $L>2$ is orthogonal under the uniform input distribution to every
such polynomial. A polynomial with the same sign as parity on every input would have positive
inner product with it. A fixed tie rule can be
absorbed into a small constant shift on the finite input set. This
restriction applies only to the single linear layer. The two-layer
variant with a gated linear unit reaches mean accuracy $0.72$ (best seed
$0.90$) at the training length and falls to chance by $L=20$.

\begin{table}[H]
\centering
\caption{Phase and readout ablations on parity. Entries give greedy
accuracy at $L{=}10$ under a matched trainer. The diagonal complex
state-space model (LRU) averages five seeds (Appendix~\ref{app:lru}).}
\label{tab:mech}
\small
\setlength{\tabcolsep}{4pt}
\begin{tabular}{lcc}
\toprule
Model & Parameters & Parity accuracy \\
\midrule
\multicolumn{3}{l}{\emph{Prior architectures and reductions}} \\
GRU (linear readout) & 1.1k & 0.50 \\
GRU + general quadratic readout & 1.6k & 0.50 \\
LSTM & 1.5k & 0.50 \\
uRNN (linear readout) & 0.8k & 0.50 \\
uRNN + squared-amplitude readout & 0.8k & 0.70 \\
LRU (linear readout) & 0.3k & 0.50 \\
LRU + squared-amplitude readout & 0.3k & 0.64 \\
Transformer (causal) & 20k & \textbf{0.99} \\
\midrule
\multicolumn{3}{l}{\emph{Our recurrence, different readouts}} \\
linear (softmax) readout & 0.4k & 0.50 \\
general quadratic readout & 2.4k & 0.58 \\
Squared-amplitude readout, phase $\theta\equiv 0$ & 0.3k & 0.56 \\
Squared-amplitude readout, with mixing $M$ & 0.3k & 0.72 \\
\wsrow \textbf{Squared-amplitude readout ($M{=}I$)} & \textbf{337} & \textbf{1.00} \\
\bottomrule
\end{tabular}
\end{table}


\begin{table}[tbp]
\centering
\caption{Diagonal complex state-space baselines on parity. All models train
at $L=10$. Entries give the five-seed mean and, in parentheses, the best seed.
A rule that predicts odd only at count five reaches $0.746$ at $L=10$.}
\label{tab:lru}
\small
\setlength{\tabcolsep}{2.5pt}
\begin{tabular}{lcccc}
\toprule
model & params & $L{=}10$ & $L{=}20$ & $L{=}80$ \\
\midrule
LRU + linear readout & 259 & 0.50 & --- & --- \\
LRU + general quadratic & 2.2k & 0.58 & --- & --- \\
LRU + squared-amplitude readout & 257 & 0.64 (0.76) & 0.50 & 0.50 \\
LRU, $|\lambda|{=}1$, + squared-amplitude readout & 257 & 0.66 (0.75) & 0.53 & 0.50 \\
LRU $\times2$ + GLU + squared-amplitude readout & 1.9k & 0.72 (0.90) & 0.50 & 0.50 \\
\midrule
\wsrow \textbf{+ observation-conditioned} & & & & \\
\wsrow \textbf{\quad rotation (Eq.~\ref{eq:evolve})} & \textbf{337} & \textbf{0.996} (1.00) & 0.57 & 0.52 \\
\bottomrule
\end{tabular}
\end{table}

\section{Symbolic Valve Task and Provenance}
\label{app:robot}
In the $\Z_4$ operate task, the policy sees the initial pointer for two
steps, the target, and the command and execution cues for intervening
turns. It learns corrective turns and the finish action from terminal
reward. Training uses 1--8 turns. The wider evaluation
includes 32, 64, and 128 turns and changes the clockwise probability.
The programmed counter is an oracle. The learned cyclic-belief counter
estimates shifts from observations and learns an action head.

The RLBench/CoppeliaSim visualization updates the valve pose with
\texttt{set\_orientation} and moves the arm with a scripted macro. Its checks
compare rendered poses with internal simulator state.
